\documentclass{article}
\usepackage{graphicx} % Required for inserting images
\usepackage[left=3cm, right=3cm, top=2cm, bottom=3cm]{geometry} %Zur Einstellung der Seitenränder
\usepackage{amsmath}
\usepackage{amssymb}
\usepackage{mathtools}
\usepackage{xcolor}
\usepackage[colorlinks=true, linkcolor=blue, citecolor=blue, urlcolor=blue]{hyperref}
\usepackage{subfiles}
\usepackage{bbold}
\usepackage{titlesec} % Manipulation der section header
\usepackage{enumitem}
\usepackage{physics}
\usepackage{empheq}
\usepackage{cancel}

%verhindert Einrücken am Beginn der nächsten Zeile nach \\
\setlength{\parindent}{0em}

% command zum Erstellen des Titels: n. Hausaufgabe AL10
\newcommand{\header}[0]{
\begin{center}
{\huge\textbf{(title)}}
\end{center}}

%define command for problem headers
\newcommand{\trieq}[0]{\overset{\wedge}{=}}

\newcommand{\uuline}[1]{\underline{\underline{#1}}}

% originally this commmand did something else, but now its just a bit obsolete, I could delete this if i just exchanged \tvec for \tilde everywhere
\newcommand{\tvec}[1]{\tilde{#1}}

\newcommand{\dist}[0]{f_{\vec{k}}}

\newcommand{\kv}[0]{\vec{k}}

\addtocontents{toc}{\protect\thispagestyle{empty}}

\begin{document}

\textcolor{red}{rename this section to something like model?}

\subsubsection{Construction of the Hamiltonian}

The Hamiltonian is 

\begin{equation}
	\mathcal{H} = \sum_{\langle ij \rangle} J_{ij}\left(\vec{S}_i\cdot \vec{S}_j + \kappa \left(\vec{S}_i\cdot \vec{S}_j\right)^2\right)
\end{equation}

Using the Abrikosov representation of the spin in terms of fermions:

\begin{equation}
	\vec{S}_i = \sum_{\alpha \beta} f_{i\alpha}^\dagger \frac{\vec{\sigma}}{2} f_{i\beta}
\end{equation}

\begin{align}
	\Rightarrow \vec{S}_i\cdot \vec{S}_j &= \frac14 \sum\limits_{\alpha \beta \gamma \delta} f_{i\alpha}^\dagger f_{i\beta} f_{j\gamma}^\dagger f_{j\delta} \underbrace{\vec{\sigma}_{\alpha \beta} \vec{\sigma}_{\gamma \delta}}_{= \delta_{\alpha \delta}\delta_{\beta \gamma} - \delta_{\alpha\beta}\delta_{\gamma \delta}} \notag \\
	&= \frac14 \sum\limits_{\alpha \beta} 2 f_{i\alpha}^\dagger f_{i\beta}f_{j\beta}^\dagger f_{j\alpha} - f_{i\alpha}^\dagger f_{i\alpha} f_{j\beta}^\dagger f_{j\beta} \notag 
\end{align}

Move the last annihilation operator in the first term to the second postion: use $ f_{j\beta}^\dagger f_{j\alpha} = \delta_{\alpha \beta} - f_{j\alpha}f_{j\beta}^\dagger$ and acquire 1 minus signs. After that bring the other annihilation operator to the last position which gives another minus.

\begin{align}
	\vec{S}_i\cdot \vec{S}_j &= \frac14 \sum\limits_{\alpha \beta} -2f_{i\alpha}^\dagger f_{j\alpha}f_{j\beta}^\dagger f_{j\beta} + 2 f_{i\alpha}f_{i\beta}\delta_{\alpha \beta} - f_{i\alpha}^\dagger f_{i\alpha} f_{j\beta}^\dagger f_{j\beta} 
\end{align}


The second term is $n_i$ and the third is $n_i n_j$. For the first term I introduce the shorthand

\begin{equation}
	\chi_{ij} = \frac12\sum_{\alpha} f_{i\alpha}^\dagger f_{j\alpha}
\end{equation}

\begin{equation}
	\vec{S}_i\cdot \vec{S}_j = -2\chi_{ij}\chi_{ji} + \frac12 n_i -\frac14 n_i n_j
\end{equation}

Using the constaint $n_i = n_j = 1$, the final expression is 

\begin{equation}
	\vec{S}_i\cdot \vec{S}_j = -2\chi_{ij}\chi_{ji} + \frac14
\end{equation}

Relevant for the Hamiltonian is also the term:

\begin{gather}
	\Rightarrow \left( \vec{S}_i \cdot \vec{S}_j \right)^2 =  \frac{1}{16}  - \chi_{ij} \chi_{ji} + 4\left(\chi_{ij}\chi_{ji}\right)^2
\end{gather}

The mean field decoupling works like 

\begin{align}
	\chi_{ij}\chi_{ji} &\approx \langle \chi_{ij} \rangle \chi_{ji} + \langle \chi_{ji} \rangle \chi_{ij} - \abs{\langle \chi_{ij}\rangle}^2 \\
	\left(\chi_{ij}\chi_{ji}\right)^2 &\approx 2\abs{\langle \chi_{ij}\rangle}^2\langle \chi_{ij} \rangle \chi_{ji} + 2\abs{ \langle \chi_{ij} \rangle}^2\langle \chi_{ji} \rangle \chi_{ij} - 3 \abs{\langle \chi_{ij}\rangle}^4
\end{align}


\begin{gather}
	\Rightarrow \vec{S}_i\cdot \vec{S}_j \approx -2\langle \chi_{ij} \rangle \chi_{ji} + \text{h.c.} + 2\abs{\langle \chi_{ij}\rangle}^2 + \frac14\\
	\Rightarrow \left( \vec{S}_i \cdot \vec{S}_j \right)^2 \approx -\langle \chi_{ij} \rangle \chi_{ji} + 8 \abs{\langle \chi_{ij} \rangle}^2\langle \chi_{ij} \rangle \chi_{ji} + \text{h.c.} + \abs{\langle \chi_{ij} \rangle}^2 - 12\abs{\langle \chi_{ij} \rangle}^4 + \frac{1}{16}
\end{gather}


\textcolor{red}{Maybe comment of the fact that the mean-field parameters know about the fermi sea and thereby about many states. In this sense the states do speak with each other}. Bringing everything together: 


\begin{align}
	\begin{split}
		\mathcal{H} \approx \sum_{\langle ij \rangle} J_{ij}&\bigg[ \left(-1 - \kappa/2 + 4 \kappa \abs{\langle \chi_{ij}\rangle}^2\right)\langle \chi_{ij} \rangle \left(f_{j\uparrow}^\dagger f_{i\uparrow} + f_{j\downarrow}^\dagger f_{i\downarrow}\right) + \text{h.c} \\
		&+ \left( 2 + \kappa - 12\kappa \abs{\langle \chi_{ij} \rangle}^2 \right) \abs{\langle \chi_{ij} \rangle}^2 + \frac{1}{4} + \frac{\kappa}{16}\bigg] 
	\end{split}
\end{align}

It is sometimes handy to define 

\begin{equation}
	\tilde{J}_{ij} = J_{ij}\left( 1+ \kappa/2\right), \quad \tilde{\kappa} = \kappa/(1+\kappa/2)
\end{equation}

I drop the constant:

\begin{align}
	\begin{split}
		\mathcal{H} \approx \sum_{\langle ij \rangle} \tilde{J}_{ij}&\bigg[ -\left(1 - 4 \tilde{\kappa} \abs{\langle \chi_{ij}\rangle}^2\right)\langle \chi_{ij} \rangle \left(f_{j\uparrow}^\dagger f_{i\uparrow} + f_{j\downarrow}^\dagger f_{i\downarrow}\right) + \text{h.c} \\
		&+ 2\left( 1 - 6\tilde{\kappa} \abs{\langle \chi_{ij} \rangle}^2 \right) \abs{\langle \chi_{ij} \rangle}^2 \bigg] 
	\end{split}
\end{align}

Hence, the mean field produces the effective hopping amplitude $t^{\text{MF}}$:

\begin{equation}
	t^{\text{MF}}_{ji} = - \tilde{J}_{ij}\langle \chi_{ij} \rangle \left(1 - 4\tilde{\kappa} \abs{ \langle \chi_{ij} \rangle}^2\right) =  -J_{ij}\langle \chi_{ij} \rangle \left(1 + \kappa/2  - 4 \kappa \abs{\langle \chi_{ij} \rangle}^2 \right).
\end{equation}

\begin{equation}
	\Rightarrow \mathcal{H} \approx \sum_{\langle ij \rangle} t^{\text{MF}}_{ji}\left( f_{j\uparrow}^\dagger f_{i\uparrow} + f_{j\downarrow}^\dagger f_{i\downarrow} \right) + \text{const.}
\end{equation}

This is pretty interesting: the hopping-amplitude $i \rightarrow j$ is directly proportional to the hopping-amplitude $j \rightarrow i$. \textcolor{red}{Here are a few ideas how to make use/discuss this insight:
\begin{itemize}
	\item Maybe one can also make the argument that we start with a "symmetric" model, i.e. a model wohtout internal sense of direction, a feature that is now preserved in the Mean field Hamiltonian. 
	\item Or maybe, even more suffisticated, one could make argument using straucture of the spin-singlet states.
	\item Maybe one can use this fact to siplify the Mean-field equation?
\end{itemize}}

\subsubsection{Continuity equation and probability current}

The probability (/particle) density on site $k$ is 

\begin{equation}
	n_k = \sum_\sigma f_{k\sigma}^\dagger f_{k\sigma}.
\end{equation}

Now I derive the corresponding probability current density:

\begin{equation}
	\frac{\partial n_k}{\partial t} = -i\sum_\sigma\left([f_{k\sigma}^\dagger, H]f_{k\sigma} + f_{k\sigma}^\dagger [f_{k\sigma},H]\right)
\end{equation}

\begin{align}
	\begin{split}
		[f_{k\sigma}^\dagger, H]f_{k\sigma} &= - \sum\limits_{\delta} t_{k+\delta;k}f_{k+\delta;\sigma}^\dagger f_{k\sigma} - \mu_k f_{k\sigma}^\dagger f_{k\sigma}\\
		f_{k\sigma}^\dagger [f_{k\sigma}, H] &= \sum\limits_{\delta} t_{k;k+\delta}f_{k\sigma}^\dagger f_{k+\delta;\sigma} + \mu f_{k\sigma}^\dagger f_{k\sigma}
	\end{split}
\end{align}

\begin{equation}
	\Rightarrow \frac{\partial n_k}{\partial t}= -\sum_\delta i \left( \sum_{\sigma} t_{k,k+\delta} f_{k\sigma}^\dagger f_{k+\delta;\sigma}-t_{k+\delta,k} f_{k+\delta;\sigma}^\dagger f_{k\sigma} \right)
\end{equation}

Define 

\begin{equation}
	J_{k \rightarrow k+\delta} \equiv i \left( \sum_{\sigma} t_{k,k+\delta} f_{k\sigma}^\dagger f_{k+\delta;\sigma}-t_{k+\delta,k} f_{k+\delta;\sigma}^\dagger f_{k\sigma} \right) 
\end{equation}

$J_{k \rightarrow k+\delta}$ is the (directed!) current density along the bond $\vec{R}_k \rightarrow \vec{R}_{k+\delta}$. The discrete (/lattice) continuity equation reads

\begin{equation}
	\frac{\partial n_k}{\partial t} = -\sum_{\delta} J_{k \rightarrow k+\delta}.
\end{equation}

Something interesting happens when calculating $\langle J_{k\rightarrow k+\delta} \rangle$:

\begin{align}
	\langle J_{k\rightarrow k+\delta} \rangle &= i\left(t_{k,k+\delta}\left\langle \sum_\sigma f_{k\sigma}^\dagger f_{k+\delta;\sigma}\right\rangle - t_{k+\delta,k} \left\langle \sum_\sigma f_{k+\delta;\sigma}^\dagger f_{k\sigma} \right\rangle\right) \notag \\
	&= 2i\left(t_{k,k+\delta}\langle \chi_{k,k+\delta} \rangle - t_{k+\delta,k} \langle \chi_{k+\delta,k} \rangle\right) \notag \\
	&=-2i \tilde{J}_{k;k+\delta}\left(1-4\tilde{\kappa}\abs{\langle\chi_{k,k+\delta}\rangle}^2\right)\bigg(\langle\chi_{k+\delta;k}\rangle\langle\chi_{k;k+\delta}\rangle - \langle\chi_{k;k+\delta}\rangle\langle \chi_{k+\delta;k}\rangle\bigg) \notag \\
	&= 0.
\end{align}

Apparently, the current on each bond vanishes in the ground state by construction. \textcolor{red}{This still leaves the possibilitie that single states have non-zero probability currents. However, it remains to investigate how this is possible for e.g. states with a singular (topologically protected) edge mode. One would guess that this feature would lead to a non-zero current in the ground state.}

\end{document}
